% SHARD-ID: AQM-09-SYMPLECTIC-SUPERCHARGE-DICTIONARY
% SHARD-TITLE: Symplectic Dictionary for CSS Supercharges
% SHARD-SUMMARY: Constructs the ghost/Koszul supercharge from an isotropic CSS stabilizer subspace.
% SHARD-SUMMARY: Proves how logical Paulis in \(L^\perp/L\) act on degree-zero supercharge cohomology.
% SHARD-SUMMARY: Separates generator-dependent and basis-free versions of the supercharge.
% SHARD-KEYWORDS: CSS code, symplectic geometry, supercharge, Koszul differential, stabilizer, logical Pauli

\section{Symplectic Dictionary for CSS Supercharges}
\label{sec:symplectic-supercharge-dictionary}

This section relates the supercharge operators to the symplectic stabilizer
geometry of Section~\ref{sec:symplectic-css-bridge}. The result is not that
the supercharge is itself a symplectic form. Rather, for any CSS stabilizer
subspace \(L\), the supercharge is the Koszul differential for the commuting
projector equations defining the \(L\)-invariant subspace.

The source anchors are registered in the local symplectic-QECC and toric-code
source manifests. The locators used here are Gottesman \path{Thesis.tex}, lines
1293--1341; Ashikhmin--Knill \path{n9.tex}, lines 249--375; and
Bombin--Martin-Delgado \path{HomologicalErrorCorrection6.tex}, lines
1206--1259.

\subsection{CSS Symplectic Data}

Work first over a prime field \(\mathbb F_p\). A CSS code on \(n\) qudits is
given by two row spaces
\[
  R_X,R_Z\subset \mathbb F_p^n
\]
with
\[
  x\cdot z=0\quad\text{for all }x\in R_X,\ z\in R_Z.
\]
The Pauli label space is
\[
  V=\mathbb F_p^n\oplus\mathbb F_p^n,
  \qquad
  \omega((x,z),(x',z'))=z\cdot x'-x\cdot z'.
\]
The CSS stabilizer label subspace is
\[
  L=(R_X\oplus 0)+(0\oplus R_Z)\subset V.
\]
The row-orthogonality condition is exactly \(L\subseteq L^\perp\).

Let
\[
  G=(\ell_1,\ldots,\ell_r)
\]
be an ordered \(\mathbb F_p\)-basis of \(L\). This basis is extra data. It is
needed to write a small generator-based supercharge, but the code subspace
depends only on \(L\).

\subsection{Projectors From Symplectic Labels}

For \(v=(x,z)\in V\), write \(W(v)=X^xZ^z\). In the whole Pauli label space,
\[
  W(v)W(w)=\zeta^{\omega(v,w)}W(w)W(v),
  \qquad \zeta=e^{2\pi i/p}.
\]
For the CSS subspace \(L=(R_X\oplus0)+(0\oplus R_Z)\), the canonical lifts are
stronger than commuting. If \(\ell=(x,z)\) and \(\ell'=(x',z')\) lie in \(L\),
then \(z\cdot x'=x\cdot z'=0\), because \(R_X\perp R_Z\). Hence
\[
  W(\ell)W(\ell')=W(\ell+\ell'),\qquad W(\ell)^*=W(-\ell).
\]
Thus \(W\) restricts to an honest representation of the additive group \(L\).

For each generator \(\ell_i\), define the stabilizer-line averaging projector
\[
  \Pi_i=\frac1p\sum_{a\in\mathbb F_p}W(a\ell_i).
\]
This is the projector onto states fixed by the one-parameter stabilizer
subgroup \(\{W(a\ell_i):a\in\mathbb F_p\}\). To check this directly,
\[
\begin{aligned}
  \Pi_i^2
  &=\frac1{p^2}\sum_{a,b}W((a+b)\ell_i)     \\
  &=\frac1p\sum_c W(c\ell_i)=\Pi_i,
\end{aligned}
\]
since for each \(c\) there are \(p\) pairs \((a,b)\) with \(a+b=c\). Also
\(\Pi_i^*=\Pi_i\), because \(W(a\ell_i)^*=W(-a\ell_i)\). Set
\[
  P_i=I-\Pi_i .
\]
Then \(P_i\) is a self-adjoint projector onto the subspace violating the
\(\ell_i\)-stabilizer equation. The projectors \(P_i\) commute with each other
because the line averages commute.

For qubits, \(p=2\), this reduces to the previous formula:
\[
  \Pi_i=\frac{I+W(\ell_i)}2,\qquad
  P_i=\frac{I-W(\ell_i)}2.
\]

\subsection{The Generator-Based Supercharge}

Attach one fermionic ghost \(\epsilon_i\) to each chosen generator
\(\ell_i\). Exterior multiplication by \(\epsilon_i\) and contraction
\(\iota_i\) satisfy
\[
  \iota_i\epsilon_j+\epsilon_j\iota_i=\delta_{ij},\qquad
  \epsilon_i\epsilon_j+\epsilon_j\epsilon_i=0,\qquad
  \iota_i\iota_j+\iota_j\iota_i=0.
\]
They commute with all physical Pauli projectors. Define
\[
  Q_G=\sum_{i=1}^r \epsilon_iP_i,\qquad
  Q_G^*=\sum_{i=1}^r \iota_iP_i .
\]
Now compute:
\[
\begin{aligned}
  Q_G^2
  &=\sum_{i,j}\epsilon_i\epsilon_jP_iP_j        \\
  &=\sum_{i<j}(\epsilon_i\epsilon_j+\epsilon_j\epsilon_i)P_iP_j=0.
\end{aligned}
\]
The anticommutator is
\[
\begin{aligned}
  Q_GQ_G^*+Q_G^*Q_G
  &=\sum_{i,j}(\epsilon_i\iota_j+\iota_j\epsilon_i)P_iP_j \\
  &=\sum_iP_i^2=\sum_iP_i .
\end{aligned}
\]
Thus the supersymmetric Hamiltonian is the stabilizer-violation penalty for
the chosen generator basis.

Degree zero means no ghost. A physical vector \(\psi\) is degree-zero closed
iff
\[
  Q_G\psi=\sum_i\epsilon_iP_i\psi=0.
\]
The one-ghost labels are independent, so this is equivalent to \(P_i\psi=0\)
for every \(i\), equivalently \(\Pi_i\psi=\psi\) for every \(i\). Since the
\(\ell_i\) form a basis of \(L\), this means
\[
  W(\ell)\psi=\psi\quad\text{for every }\ell\in L.
\]
There are no degree \(-1\) states, so no degree-zero exact states. Therefore
\[
  H^0(Q_G)
  =
  \{\psi:W(\ell)\psi=\psi\text{ for all }\ell\in L\},
\]
the CSS stabilizer code space. If \(r=\dim L\), its dimension is \(p^{n-r}\).

\subsection{Logical Paulis Commute With the Supercharge}

Let \(v\in L^\perp\). Then \(\omega(v,\ell_i)=0\) for every \(i\), so \(W(v)\)
commutes with every \(W(a\ell_i)\), hence with every \(\Pi_i\), \(P_i\), and
\[
  Q_G=\sum_i\epsilon_iP_i.
\]
Therefore \(W(v)\) acts on \(H^0(Q_G)\).

If \(v\in L\), then \(W(v)\) acts as the identity on \(H^0(Q_G)\) in the
trivial-character sector. Hence the action factors through
\[
  L^\perp/L.
\]
This is the exact point where symplectic geometry enters the supercharge
cohomology: the degree-zero cohomology carries the logical Pauli action of the
symplectic reduction.

If \(v\notin L^\perp\), then some \(\ell_i\) has \(\omega(\ell_i,v)\neq0\).
For \(\psi\in H^0(Q_G)\),
\[
  W(a\ell_i)W(v)\psi
  =
  \zeta^{a\omega(\ell_i,v)}W(v)\psi .
\]
Averaging over \(a\) gives \(\Pi_iW(v)\psi=0\). Thus
\[
  P_iW(v)\psi=W(v)\psi,
  \qquad
  Q_GW(v)\psi
  \text{ has a nonzero } \epsilon_i\text{-component}.
\]
So Pauli labels outside \(L^\perp\) are precisely those detected by the
supercharge as nonzero stabilizer syndrome.

\subsection{Basis Dependence and Basis-Free Variant}

The operator \(Q_G\) depends on the chosen basis \(G\) of \(L\). Changing
generators changes the individual projectors \(P_i\) and therefore changes the
positive Hamiltonian \(\sum_iP_i\). The degree-zero cohomology does not change:
it is always the common invariant subspace for the whole \(L\).

There is also a basis-free, but redundant, version. Let \(\mathbb P(L)\) be
the set of one-dimensional subspaces of \(L\). For a line
\(\lambda=\mathbb F_p\ell\), define
\[
  \Pi_\lambda=\frac1p\sum_{a\in\mathbb F_p}W(a\ell),
  \qquad
  P_\lambda=I-\Pi_\lambda .
\]
This is independent of the nonzero representative \(\ell\). Attach one ghost
\(\epsilon_\lambda\) to each line and set
\[
  Q_L^{\mathrm{all}}
  =
  \sum_{\lambda\in\mathbb P(L)}\epsilon_\lambda P_\lambda .
\]
The same CAR/projector proof gives \((Q_L^{\mathrm{all}})^2=0\). Its
degree-zero cohomology is again the CSS code space. The price is redundancy:
if \(\dim L=r\), then
\[
  |\mathbb P(L)|=\frac{p^r-1}{p-1}.
\]
The generator-based \(Q_G\) is smaller; the projective-line \(Q_L^{\mathrm{all}}\)
is intrinsic to \(L\).

\subsection{General CSS Dictionary}

For a CSS code with check matrices \(H_X,H_Z\) over \(\mathbb F_p\), the
dictionary is summarized by the table below. In the cellular toric-code case,
this means
\[
  Q_{\mathrm{cell}}=\delta^0+\delta^1
  \quad\leadsto\quad
  R_X=\operatorname{im}\delta^0,\quad R_Z=\operatorname{im}\partial_2
  \quad\leadsto\quad
  L\quad\leadsto\quad Q_G .
\]
The cellular differential supplies the two CSS support spaces. The symplectic
step packages them as \(L\). The ghost/Koszul step turns a chosen basis of \(L\)
into commuting projectors and the operator \(Q_G\).
\[
\begin{array}{c|c}
\text{CSS/symplectic object} & \text{supercharge object}\\
\hline
R_X,R_Z & X\text{- and }Z\text{-type stabilizer equations}\\
L=(R_X\oplus0)+(0\oplus R_Z) & \text{common invariant equations}\\
L\subseteq L^\perp & \text{commuting projectors }P_i\\
\text{basis }G\text{ of }L & \text{minimal ghost list}\\
P_i=I-\frac1p\sum_aW(a\ell_i) & \text{violated-check projector}\\
Q_G=\sum_i\epsilon_iP_i & \text{Koszul supercharge}\\
L^\perp & \text{Paulis commuting with }Q_G\\
L^\perp/L & \text{logical Paulis on }H^0(Q_G)
\end{array}
\]
Thus the supercharge is a homological-algebra package built on top of an
isotropic subspace. Symplectic geometry supplies \(L\), its normalizer
\(L^\perp\), and the logical quotient; the supercharge resolves the common
stabilizer equations defining the code space.

\subsection{Prime-Power Note}

For \(q=p^m\), replace \(\mathbb F_p\) by \(\mathbb F_q\), and use the additive
character
\[
  \chi(t)=\exp(2\pi i\,\operatorname{Tr}_{\mathbb F_q/\mathbb F_p}(t)/p).
\]
For a line \(\mathbb F_q\ell\), the invariant projector is
\[
  \Pi_\ell=\frac1q\sum_{a\in\mathbb F_q}W(a\ell).
\]
The same calculation proves \(\Pi_\ell^2=\Pi_\ell\), commuting projectors for
isotropic \(L\), \(Q^2=0\), and
\[
  H^0(Q)\cong\text{the }L\text{-stabilizer code}.
\]
For composite non-field qudit dimension, one must replace vector spaces by
modules over \(\mathbb Z_D\), and the clean count \(p^{n-\dim L}\) need not be
imported without checking torsion.

\subsection{Exact Validation}

The script \path{css_supercharge_symplectic_dictionary.jl} in
\path{scripts/lattice_codes/} validates the dictionary without building Hilbert
matrices. It checks the binary Steane CSS code and a qutrit CSS toy code. The
CSV \path{css_supercharge_symplectic_dictionary.csv} records CSS isotropy,
stabilizer rank, code dimension, logical symplectic dimension, ghost counts,
and the \(Q^2\) and anticommutator certificates in the
\path{2026-05-24-css-supercharge-symplectic-dictionary} run bundle.
